\section{Separating Hindsight Memorization from Foresight}
\label{sec:foresight}

The deepest objection to any backtest run with a modern LLM anywhere in
the loop (Section~\ref{sec:threats}) is that its apparent performance
decomposes into three terms:
\begin{equation*}
\text{Performance} \;=\;
\underbrace{\text{reasoning over pre-cutoff evidence}}_{\text{what we want}}
\;+\; \underbrace{\text{memorized future}}_{\text{contamination}}
\;+\; \underbrace{\text{topic prior}}_{\text{fashion}}
\end{equation*}
and a single retrospective instance cannot tell the terms apart. A
question that 2024 answered may have been \emph{predicted} from 2020
evidence---or \emph{remembered} from the model's training data. This
section reports two experiments designed to pry the terms apart:
holding the cutoff fixed while varying which component generates the
question, and holding the generators fixed while moving the cutoff
across the judge-model's training boundary.

\subsection{Same cutoff, different generators: where does the signal
live?}
\label{sec:foresight:abc}

Three generation pipelines share the 2020 cutoff and the v1L evaluation
conditions but differ in what produces the question:

\begin{description}
  \item[A --- LLM only.] The direct-LLM baseline: \texttt{gpt-4.1} reads
    sampled pre-cutoff abstracts and proposes questions. Weights fully
    exposed to post-2020 literature.
  \item[B --- structure $\to$ LLM.] A deterministic, LLM-free detector
    finds pairs of pre-cutoff abstracts about the same catalogued object
    with opposing stances on the same species (detection vs.\
    non-detection / upper limit); \texttt{gpt-4.1}'s only job is to
    verbalize each detected tension as one falsifiable question. The
    evidence structure is fixed before any LLM sees anything.
  \item[C --- structure only.] The same detected pairs rendered by a
    fixed template. No LLM anywhere: this generator has no weights to
    contaminate.
\end{description}

B and C share identical evidence structures, so their gap isolates the
language-realization layer; A and B share the same LLM, so their gap
isolates evidence structure. (B/C are an evidence-structure-\emph{lite}
probe---object co-mention plus stance cues---not a reimplementation of
the evidence-graph system.) Table~\ref{tab:abc} gives the $n=125$
results.

\begin{table}[t]
  \centering
  \small
  \caption{The A/B/C decomposition on astronomy v1L (judge-only,
  $n=125$ each; brackets are 95\% Clopper--Pearson intervals).
  $s_{\mathrm{obj}}$ = share of questions naming a specific catalogued
  object (deterministic rubric, Section~\ref{sec:foresight:saf}).}
  \label{tab:abc}
  \begin{tabular}{@{}lllllr@{}}
    \toprule
    Pipeline & Eng. & Ans. & Ref. & $s_1$ & $s_{\mathrm{obj}}$ \\
    \midrule
    A: LLM only              & 96\% [91,99] & 15\% [9,23]  & 0\% [0,3]   & 0.722 & 5\% \\
    B: structure $\to$ LLM   & 74\% [65,81] & 39\% [31,48] & 13\% [8,20] & 0.694 & 95\% \\
    C: structure only        & 77\% [68,84] & 25\% [17,33] & 2.4\% [0,7] & 0.721 & 100\% \\
    \bottomrule
  \end{tabular}
\end{table}

Three readings. \emph{First}, the contamination-free pipeline C beats
the fully exposed pipeline A on answered rate (24.8\% vs.\ 15.2\%,
$p=0.08$) and on refutations (3 vs.\ 0)---evidence that a foresight
signal exists in pre-cutoff evidence structure alone, extractable with
zero model weights. \emph{Second}, adding the LLM back as a pure
verbalizer (B) roughly doubles resolution over the same structure
(39.2\% vs.\ 24.8\% answered, $p=0.02$; 12.8\% vs.\ 2.4\% refuted,
$p=0.003$): phrasing a tension as a crisp either/or question makes it
judgeable, and B ends with five times pipeline A's refutation count
while citing the same model. \emph{Third}, A's questions are
structurally different, not just weaker: only 5\% name a specific
object (vs.\ 95--100\% for B/C), and its engagement is the highest of
any system---broad questions that everything touches and little
settles.

\subsection{Same generators, moving cutoff: the temporal stress test}
\label{sec:foresight:stress}

If pipeline A's performance were substantially the memorized-future
term, it should degrade as the cutoff crosses the model's training
boundary (June 2024 for \texttt{gpt-4.1}). We ran four generators (A,
B, C, and random claims as an era control) at four cutoffs---2010,
2015, 2020, 2024---with uniform four-year future windows ($n=50$ per
cell; the 2024 window is censored at 2026-06 and flagged; corpora per
era rebuilt from released manifests). Table~\ref{tab:stress} reports
the grid.

\begin{table}[t]
  \centering
  \small
  \caption{Temporal contamination stress test (judge-only, $n=50$ per
  cell; c2010 tension cells have $n=48$---the 855-paper 2005--2010
  corpus yields only 48 detectable tension pairs). Cutoffs 2010--2020
  lie inside the LLM's training data; the 2024 cutoff's future window
  (2025--2026-06, censored) postdates it. Ans.\ = answered; Ref.\ =
  refuted count.}
  \label{tab:stress}
  \begin{tabular}{@{}llcccc@{}}
    \toprule
    & & c2010 & c2015 & c2020 & c2024 \\
    \midrule
    A: LLM only & Ans. & 22\% & 12\% & 10\% & 16\% \\
                & Ref. & 2 & 0 & 0 & 0 \\
    B: structure $\to$ LLM & Ans. & 31\% & 26\% & 34\% & 32\% \\
                & Ref. & 6 & 3 & 6 & 5 \\
    C: structure only & Ans. & 10\% & 6\% & 16\% & 4\% \\
                & Ref. & 0 & 0 & 0 & 0 \\
    Random claims & Ans. & 12\% & 6\% & 4\% & 6\% \\
                & Ref. & 1 & 1 & 0 & 0 \\
    \bottomrule
  \end{tabular}
\end{table}

\paragraph{The naive collapse does not happen.} Pipeline A's answered
rate is statistically flat across the training boundary (14.7\% pooled
in-training vs.\ 16.0\% post-training, $p=0.82$), as is every other
generator's. At the outcome level, the memorized-future term is
\emph{not} where A's performance comes from---because, as the
decomposition shows, A's performance never rested on specifics that
memorization could supply. Its engagement sits at 92--98\% at every
cutoff, its specificity at the floor ($\approx$1.1 of 3) at every
cutoff: broad questions about each era's active topics, engaged
everywhere, resolving little, in any era. That is the topic-prior term
at work, and a topic prior does not need to remember the future---the
present is enough.

\paragraph{Where a memorization trace does appear.} Two places, both
in the channels the topic prior cannot supply. Pipeline A's only
premise refutations in the entire stress test (2 of 200) occur at the
deepest-contamination cutoff, 2010---the one era whose reversals the
model has certainly read about---and never after (0 of 150; too few
for significance, CI [0.5\%, 13.7\%] at c2010). And A's phrasing
similarity to future literature is highest inside its training window
(0.717--0.724) with a mild post-training dip (0.703)---directionally
consistent with phrasing-level memorization, though era confounds keep
this suggestive rather than conclusive.

\paragraph{The structural signal is era-robust.} Pipeline B refutes
premises at every cutoff---6, 3, 6, 5---including the one whose future
the model cannot have seen (10.0\% post-training vs.\ 10.1\%
in-training). Across all cutoffs B refutes at 10.1\% vs.\ A's 1.0\%
($p=4\times10^{-5}$), and B's answered margin over C persists
post-training ($+28$ points at c2024). One residual channel remains
open and is worth stating precisely: c2024 tension pairs are drawn from
2019--2024 abstracts, and a tension resolved in early-2024 literature
the model saw could steer B's phrasing even though the evaluation
window postdates training. C is immune by construction, which is why C
finding \emph{any} future engagement at every cutoff (70--83\%) is the
cleanest single fact in the grid.

\subsection{Specificity-adjusted foresight}
\label{sec:foresight:saf}

Coverage can be farmed by asking broad questions---``How can we better
understand exoplanet atmospheres?'' will be engaged with probability 1
in any active field. We therefore score every question with a
deterministic, released rubric: $+1$ for naming a specific catalogued
object, $+1$ for naming a measurable claim (species, quantity with
units, abundance comparative), $+1$ for an explicit discriminative
construction (``\ldots or an artifact of\ldots'', ``as tested by'');
and define specificity-adjusted foresight
$\mathrm{SAF} = \mathbb{E}[\,w(\text{outcome}) \cdot \text{specificity}/3\,]$
with $w$ = 1 / 0.5 / 0.25 / 0 for answered / partial / posed-open /
not addressed. (Template pipelines inherit the test-construction
point from their template---the rubric's components are reported
separately for exactly that reason; a novelty term awaits first-posed
dates, Section~\ref{sec:conclusion}.) Two facts survive every era and
instance: pipeline A's anchoring is an order of magnitude below the
structure pipelines' ($s_{\mathrm{obj}}$ 5\% vs.\ 95--100\% on v1L),
and its SAF never exceeds the random-template floor by more than a few
points (0.18--0.23 vs.\ 0.24--0.26), while structure pipelines reach
0.34--0.43. The evidence-graph submission's ten questions score
$s_{\mathrm{obj}}=90\%$, SAF 0.40.

\subsection{What this section establishes}

\begin{quote}
\emph{Memorized relevance is not scientific foresight.}
\end{quote}

The LLM-only pipeline exhibits relevance everywhere---near-ceiling
engagement, the closest phrasing to the future literature at every
cutoff---and specific foresight nowhere: no refutation outside the era
it could have memorized, an answered rate indistinguishable from random
templates, specificity at the floor. The structure-first pipelines
invert the picture: lower engagement, higher resolution, refutations in
every era including the one that postdates the model's training. On
this evidence, the foresight signal measured by historical backtesting
lives chiefly in pre-cutoff evidence structure; the LLM contributes a
real but separable service---turning a detected tension into a
question sharp enough to be judged. Contamination, meanwhile, turns
out to be measurable rather than merely confessable: the stress test
bounds its outcome-level effect (small for this task family) and
localizes its traces (phrasing proximity; refutations only in deep
history). All 798 stress-test judgments, the A/B/C submissions, and
the per-cell corpora manifests are released; every number in this
section is judge-only and carries $n=48$--$125$ intervals---the
qualitative pattern, not any single rate, is the finding.
